\documentclass[../../main.tex]{subfiles}

\begin{document}
\chapter{量子理想气体}

受限于数学手段，我们一直未能处理完全不可分辨的量子系统，对这种系统，我们不仅要考虑能量量子化，还要小心地处理全同粒子的不可分辨性，这意味着我们要区分费米子和玻色子，自然我们也要考虑粒子的自旋。根据上一章发展出的巨正则系综理论，我们有信心统计出近独立的玻色系统和费米系统的各热力学量，而这种系统的一个很恰当例子就是气体系统，我们称之为量子理想气体，细分下来又分为玻色气体和费米气体。我们先讨论两种特殊的玻色气体，接着从一般的角度讨论一下量子理想气体，最后分别针对玻色气体和费米气体，我们将讨论一些独属于它们自身的特性。

\section{再谈电磁辐射——光子气体}\label{sec:光子气体}

本节的标题是开门见山类型的，我们要再一次讨论平衡辐射的统计理论，只是本节我们不把平衡辐射视为电磁波，而是一种粒子，即光子。从热力学量预言上来说，本节不会给出任何新的内容，我们的介绍只是让读者从另一个视角理解电磁辐射，这种视角将启发我们对固体振动的另一种认识，推而广之，这是一种广泛使用的视角，对读者日后的学习总是有裨益的。

\subsection{光量子}
读者应该都知道，关于光的本性之争，即光的波动说和微粒说，最后被光的波粒二象性\footnote{请读者谨慎对待波粒二象性，光的波粒二象性不是说光即使波又是粒子，而是光既有波的某些属性也有粒子某些属性。光的波动属性是光有频率和波长，可以发生干涉衍射等；光的粒子属性是光可以被计数，我们能讨论一个光子或者多个光子。从本性来看，光既不是经典的电磁波也不是经典的粒子，光子的全称是光量子，是一种量子性质的粒子。类似的实物粒子如电子的波粒二象性也表明电子既不是经典的粒子也不是经典的波，当我们说电子是一个波的时候，指的是概率波而不是经典的机械波或电磁波。可以看出波粒二象性是一个否定性的本体认识，现代物理到底怎么任何光和电子等物质，读者要去量子场论中学习。}所终结，人们意识到必须同时承认光的波动性和粒子性。\ref{sec:电磁辐射的统计理论}节我们用正则系综处理电磁辐射，统计的起点是电磁场能量的量子化即电磁场的能量是谐振子的能量和，这已经暗示了电磁场的粒子性，但我们直接把平衡辐射视为谐振子系统，用正则系综给出了统计结果。这种视角很便于统计，但这也抹去了我们对平衡辐射的粒子性的探索。下面我们换种视角，把平衡辐射看成所谓的光子系统。

平衡辐射的能量是谐振子的能量和，而每个谐振子的频率是简正模式的频率，这一观点的数学表述是
\begin{equation}\label{eq:平衡辐射谐振子}
E=\sum_{\bm{k},\lambda}\left(n_{\bm{k},\lambda}+\frac{1}{2}\right)\hbar\omega_{\bm{k},\lambda}\quad \bm{k}=\frac{\pi}{L}(m,n,p),\lambda=\pm1
\end{equation}
上式中\(\bm{k}\)对应简正模式的三个本征值，而\(\lambda\)表示两个独立的偏振方向\footnote{\(\lambda=\pm1\)完全是人为规定，显然它的值不影响简正模式的频率，只要它有两个值就行}，这四个值合在一起确定了一个简正模式。我们略去零点能，现在关注某个特定的简正模式，它对应的能量为\(n_{\bm{k},\lambda}\hbar\omega_{\bm{k},\lambda}\)，不难看出这个能量是\(\hbar\omega_{\bm{k},\lambda}\)的\(n_{\bm{k},\lambda}\)倍。如果我们假想一种粒子，它的能量为\(\hbar\omega_{\bm{k},\lambda}\)，显然\(n_{\bm{k},\lambda}\)就是这种粒子的个数。我们现在考虑的体系有无数个简正模式，显然不可能假想出无数种粒子，我们只需规定刚才假想的粒子，它有无数个量子态，每个简正模式\((\bm{k},\lambda)\)对应一个量子态，而这个量子态的能量为\(E_{\bm{k},\lambda}=\hbar\omega_{\bm{k},\lambda}\)。根据以往的学习经验，读者不难看出，这里假想的粒子正是光量子，简称光子。爱因斯坦提出光量子假说时，认为光的能量是量子化的，把光的能量的基本单元称为光子，这里的思路也是这样的，只不过爱因斯坦作为理论的探索者，我们无从得知它的假说从何而来，我们这里的出发点是式\eqref{eq:平衡辐射谐振子}，这是量子场论的结论\footnote{实际上本段引入光子的思路就是量子场论中引入光子的思路}。

引入光子后，我们对式\eqref{eq:平衡辐射谐振子}的看法就不是谐振子的能量和了，而是光子的能量和。其中\(\hbar\omega_{\bm{k},\lambda}\)表示光子某一量子态的能量，\(\frac{1}{2}\)表示零点能，后文不作考虑。\(n_{\bm{k},\lambda}\)表示这个量子态上的光子数，它本质上是表示谐振子能级的量子数，因此它可以是任意非负整数，这意味这光子不受泡利不相容原理的限制，它是玻色子。光子的存在是从能量量子化中抽象出来的，在我们的介绍中，光子是假想出的粒子，因此它和前文涉及的非相对论或极端相对论自由粒子有很多不同。比如光子没有如自由实物粒子一样的波函数，讨论光子的位置完全没有意义。但这对统计没有影响，我们多次提及，统计物理并不在意粒子的具体状态，更重要的是状态的能量和计数\footnote{我们的课程没有介绍谐振子的定态波函数，但不妨碍我们在统计中广泛使用谐振子模型}。不过把光子的能级公式完整写出
\begin{equation}
    E_{\bm{k},\lambda}=\hbar\omega_{\bm{k},\lambda}=\frac{\hbar c\pi}{L}\sqrt{m^2+n^2+p^2}
\end{equation}
能看出来这和极端相对论自由粒子的能级公式很相似。其中原因我们在推导极端相对论自由粒子的能级公式中已经透露，自由粒子的定态的离散和经典波的简正模式的离散相同，而光子的量子态就来自于电磁波的简正模式。但是细致考虑起来，两者有一个很大的不同，对于实物粒子来说，能量动量关系是已知的，通过德布罗意关系式\(\bm{p}=\hbar\bm{k}\)求出了能级公式。而光子的能级公式是直接从量子化的电磁场能量中抽象出来的，根据电磁波的色散关系\(\omega=ck\)，我们能得到光子的能量动量关系，即
\begin{equation}
    E=\hbar\omega=\hbar ck=pc
\end{equation}
上式\footnote{可能在读者以往的学习经历中，光子速度为光速质量为0，满足\(E=pc\)是前提条件，基于此和\(E=\hbar\omega=\hbar ck\)得到了\(p=\hbar k\)，在成熟的量子理论中，这一逻辑有问题，因为光子不再只是爱因斯坦的奇思异想，它的引入有严格的数理基础，对光子的各项性质都不能想当然}的成立没有近似条件，根据\(E^2=p^2c^2+m^2c^4\)，这表明光子的质量为0，或者说光子的速度为光速。

光子作为能量的基本单元使得它和实物粒子还有一个显著的区别，即光子数不守恒。为了理解光子数不守恒，我们先讨论一下实物粒子系统。考虑一个由同种粒子构成的封闭系统，直觉告诉我们系统中的粒子总数是固定的，不会多也不会少。进一步考虑，这种直觉来自两点前提：一是系统的总质量不会变，二是同种粒子的质量相同。当然，如果系统中的粒子带电，那么粒子数固定也可以来自电荷守恒和每种粒子的电荷相同。也就是说，当我们谈及某个粒子的粒子数时，一定意味着我们选取了该粒子的某个基本的不变的属性，而且这个属性是可以累加的，对于一群粒子，通过这群粒子的属性总值除以单个粒子的属性值，才能数出这群粒子到底有几个。而这群粒子的个数不变意味这这个属性的总值是不变的。对于光子，我们找不出这样的一种属性，所以我们不能直接数出一个系统的光子数。实际上前文谈及光子的粒子数指的都是某一量子态的光子数，它可以通过该量子态占据的能量和该量子态的能级数出来，系统的总光子数由各量子态的光子数加和而来。正因为总光子数的这一特点，系统的光子数是不守恒的，虽说不是任意的，但绝对不是一个定值。数学化地说就是，假设光子系统的总能量固定即\(E=\sum_{\bm{k},\lambda}n_{\bm{k},\lambda}E_{\bm{k},\lambda}\)不变，那么光子总数\footnote{再次强调，对于实物粒子\(N=\sum_r n_r\)是对各量子态粒子数的约束，对于光子\(N=\sum_r n_r\)就是光子总数的定义}\(N=\sum_{\bm{k},\lambda}n_{\bm{k},\lambda}\)没有任何理由不变，系统的光子数是不固定的，或者准确的说是不守恒的\footnote{对于开放的实物粒子系统，虽然它的粒子数不固定，但连上环境总的粒子数是固定的，所以光子数不仅仅是不固定，连上光子的环境即空窖壁，光子数仍然不固定}。倘若读者觉得这些分析过于繁琐，考虑如下的简单情形也能告诉我们光子数不守恒的事实。对于一个实物粒子构成的开放系统，倘若系统的粒子数多一，那么按常理一定是环境的粒子数少一，一个粒子从环境进入系统。但对于空窖中的平衡辐射，假设系统中多了一个光子，那么需要环境即空窖壁中减少一个光子吗？并不需要，我们只需要让空窖壁中的某个原子跃迁至低能级，那么空窖壁就向平衡辐射中辐射了一个光子，显然空窖壁辐射出一个光子不会使空窖壁的光子数减少，因为空窖壁中根本没有光子！

我们总结一下我们对光子的认识，尤其是统计上的认识。光子显然是相互独立的，因为光子仅仅作为能量单位被抽象出来，完全不涉及相互作用能。光子作为构成系统的粒子，它们之间是独立的，类似我们之前处理的盒子中的自由粒子，而这正是我们对理想气体系统的建模\footnote{我们处理的都是无相互作用系统，但理想气体特殊在它明摆着就是无相互作用的，而不像固体热容理论，各粒子不独立但简正模式独立}，所以我们不妨把平衡辐射称为光子气体。对于近独立的光子气体系统，我们已经知道了光子的能级，还需要知道光子是否可分辨。同样因为光子是能量单位，某个量子态有若干个光子指的是若干份能量，我们不可能去分辨能量，所以光子是不可分辨的。再结合前文光子是玻色子的结论，我们说，光子气体是玻色系统。按照之前的学习经验，我们要先分析一下系统是否满足非简并条件，这意味着我们可以沿用之前的统计手段。但对于光子，粒子数不守恒带来了一个重要结论，即光子气体的化学势为0，这为我们使用巨正则系综计算玻色系统的热力学量提供了便利，我们无需在意光子气体是否满足非简并条件。下面我们分析为什么光子气体的化学势为0。

从不严谨的角度说，开放系统的粒子数变化意味着环境的粒子数变化，而光子系统的粒子数变化不需要外界发生任何改变，仅仅通过再分配各量子态所占据的能量就能改变光子数，改变光子数没有代价，所以光子气体的化学势为0。严谨的数学计算同样从开放系统说起。对于一个开放系统，设其为\(A\)同时设环境为\(B\)，两者共同构成了孤立系统，我们仅考查开放系统粒子数变化会导致的结果。孤立系统平衡时熵\footnote{根据\ref{subsec:宏观系统的最概然特性}小节最后的论述，这里的熵指的是系统在某种分配下的熵}取极值意味着熵对任何自由变量的导数为0，我们选取系统\(A\)的粒子数为自由变量，那么\begin{equation}
    \pdv{S}{N_A}=\pdv{S_A}{N_A}+\pdv{S_B}{N_A}=0
\end{equation}
考虑到开放系统和环境的总粒子数固定，所以\(\dd N_A=-\dd N_B\)即\(\pdv{N_B}{N_A}=-1\)，上式移项得
\begin{equation}
    \pdv{S_A}{N_A}=-\pdv{S_B}{N_A}=-\pdv{S_B}{N_B}\pdv{N_B}{N_A}=\pdv{S_B}{N_B}
\end{equation}
有热力学基本微分方程不能看出上式意味着
\begin{equation}
    \mu_A=\mu_B
\end{equation}
这是热力学中就有的结论，前文我们说开放系统与环境交换粒子最终两者化学势相同就来自于此。对于光子系统，我们做类似的考虑，但光子数不守恒意味着光子系统的粒子数变化不会导致环境粒子数的改变即
\begin{equation}
    \pdv{N_B}{N_A}=0
\end{equation}
所以
\begin{equation}
    \pdv{S_A}{N_A}=-\pdv{S_B}{N_A}=-\pdv{S_B}{N_B}\pdv{N_B}{N_A}=0
\end{equation}
显然这意味\(\mu_A\)即光子气体的化学势为0。

\subsection{光子气体的统计}
上一小节，我们引入了光子的概念，阐述了对光子的统计物理认识，这些认识足够我们开始对光子气体进行统计。玻色系统的巨配分函数为
\begin{equation}
    \ln\mathcal{Z}=-\sum_r \ln(1-\e^{-\alpha-\beta E_r})
\end{equation}
光子气体的化学势为0即\(\alpha=0\)，带入上式会使巨配分函数的表达式更简洁，但这也会导致巨配分函数中不再出现\(\alpha\)以至于我们无法计算系统的粒子数。所以我们先保留\(\alpha\)，得到热力学量的表达式后再令\(\alpha=0\)。对于光子气体，巨配分函数中的求和也无法解析计算，读者应该猜到怎么做了，我们要用态密度，把求和变成积分。课程进行到这，读者对自由粒子态密度的计算应该是轻车熟路的，这里给出计算过程但不再做任何说明。关于波矢大小\(k\)的态密度为
\begin{equation}
    g(k)\dd k=2\frac{\frac{1}{8}4\pi k^2\dd k}{\left(\frac{\pi}{L}\right)^3}=\frac{Vk^2}{\pi^2}\dd k
\end{equation}
由\(E=\hbar\omega=\hbar ck\)可得\footnote{利用\(\omega=ck\)可得关于\(\omega\)的态密度，读者不难发现这里计算的一系列态密度和\ref{sec:电磁辐射的统计理论}节中的简正模式密度完全一致，因为我们把每个简正模式抽象成了一个量子态}
\begin{equation}
    g(E)\dd E=g(k)\dd k=\frac{V}{\pi^2(\hbar c)^3}E^2\dd E
\end{equation}
这给出巨配分函数
\begin{equation}
    \begin{split}
        \ln\mathcal{Z}&=-\frac{V}{\pi^2(\hbar c)^3}\int_0^{\infty}E^2\ln(1-\e^{-\alpha-\beta E})\dd E \\
        &=-\frac{V}{3\pi^2(\hbar c)^3}\int_0^{\infty}\ln(1-\e^{-\alpha-\beta E})\dd E^3 \\
        &=-\frac{V}{3\pi^2(\hbar c)^3}\left[\eval{E^3\ln(1-\e^{-\alpha-\beta E})}_{0}^{\infty}-\int_0^{\infty}\frac{E^3\beta}{\e^{\alpha+\beta E}-1}\dd E\right] \\
        &=\frac{V\beta}{3\pi^2(\hbar c)^3}\int_0^{\infty}\frac{E^3}{\e^{\alpha+\beta E}-1}\dd E
    \end{split}
\end{equation}
分部积分产生的常数项为0，请读者自行验证。最后的积分由于存在\(\alpha\)，无法积出初等函数解\footnote{后文我们会再次遇到类似的积分，到时会把这个积分写成一个特殊函数，可以说本章所有的数学计算难点都在于积分中的这一个\(\alpha\)。倘若此处的积分没有\(\alpha\)，虽说仍没有初等解析解，但这种积分结果很简洁，正是玻色积分}，我们先保留该积分不做计算。求出热力学量后令\(\alpha=0\)再计算新的积分。

由巨配分函数得，光子气的内能为
\begin{equation}
    \begin{split}
        U=-\pdv{\ln\mathcal{Z}}{\beta}=&\frac{V\beta}{3\pi^2(\hbar c)^3}\int_0^{\infty}\frac{E^4\e^{\alpha+\beta E}}{(\e^{\alpha+\beta E}-1)^2}\dd E \\
        \xlongequal{\alpha=0}&\frac{V\beta}{3\pi^2(\hbar c)^3}\int_0^{\infty}\frac{E^4\e^{\beta E}}{(\e^{\beta E}-1)^2}\dd E \\
        \xlongequal{x=\beta E}&\frac{V}{3\pi^2(\hbar c)^3\beta^4}\int_0^{\infty}\frac{x^4\e^{x}}{(\e^{x}-1)^2}\dd E
    \end{split}
\end{equation}
在得到热力学量后我们令\(\alpha=0\)，这给出一个可计算的积分，我们将看到，它能化简为玻色积分。
\begin{equation}
    \begin{split}
        \int_0^{\infty}\frac{x^n\e^{x}}{(\e^{x}-1)^2}\dd E&=-\int_0^{\infty}x^n\dd{\frac{1}{\e^x-1}} \\
        &=-\eval{\frac{x^n}{\e^x-1}}_{0}^{\infty}+n\int_0^{\infty}\frac{x^{n-1}}{\e^x-1}\dd x \\
        &=n\zeta(n)\Gamma(n)=\zeta(n)n!
    \end{split}
\end{equation}
上式的运算默认了\(n\neq1\)，这样才能保证分部积分产生的常数为0，同时也保证了玻色积分本身不发散\footnote{读者可用换元法计算\(\int_0^{\infty}\frac{1}{\e^x-1}\dd x\)，它确实是发散的，玻色积分\(\int_0^{\infty}\frac{x^n}{\mathrm{e}^x-1}\dd x=\zeta(n+1)\Gamma(n+1)\)中\(\zeta(1)\)为调和级数，也说明了积分发散}。根据上式，\(\int_0^{\infty}\frac{x^4\e^{x}}{(\e^{x}-1)^2}\dd E=\zeta(4)\times 4! =\frac{\pi^4}{90}\times24=\frac{4\pi^4}{15}\)，这给出光子气体的内能密度
\begin{equation}
    u=\frac{4\pi^2\k^4}{45\hbar^3c^3}T^4
\end{equation}
\end{document}